\chapter{Tests, Validation et Déploiement}

% ============================================================
\section{Introduction}
% ============================================================

Ce dernier chapitre technique présente la stratégie de tests adoptée pour valider
les 34 user stories du projet, la chaîne DevOps (Docker, CI/CD), l'analyse de
qualité logicielle via SonarCloud, et l'infrastructure de déploiement retenue.

% ============================================================
\section{Tests fonctionnels}
% ============================================================

Les tests fonctionnels ont été conduits en suivant les scénarios nominaux et
alternatifs de chaque user story. Ils couvrent les six sprints et vérifient le
comportement attendu de l'API depuis le point de vue de l'appelant.

\begin{table}[H]
  \centering
  \caption{Résultats des tests fonctionnels par sprint}
  \label{tab:tests_fonc}
  \begin{tabularx}{\textwidth}{L{5cm}C{2cm}C{2cm}L{3cm}}
    \toprule
    \thead{Module} & \thead{Scénarios} & \thead{Réussis} & \thead{Couverture} \\
    \midrule
    \rowcolor{lightgray}
    Auth. \& utilisateurs    & 12 & 12 & 100\,\% \\
    Projets, équipes, sprints & 15 & 15 & 100\,\% \\
    \rowcolor{lightgray}
    Tâches, Kanban, Gantt    & 18 & 18 & 100\,\% \\
    Notifications, calendrier & 10 & 10 & 100\,\% \\
    \rowcolor{lightgray}
    IA (ML + LLM)            &  9 &  9 & 100\,\% \\
    RBAC \& journal d'audit   & 14 & 14 & 100\,\% \\
    \midrule
    \textbf{Total} & \textbf{78} & \textbf{78} & \textbf{100\,\%} \\
    \bottomrule
  \end{tabularx}
\end{table}

Les tests d'API ont été réalisés avec \textbf{Postman} (collections organisées
par sprint) puis exécutés automatiquement avec \textbf{Newman} pour les scenarii
de régression.

% ============================================================
\section{Assurance qualité}
% ============================================================

\subsection{SonarCloud}

\textbf{SonarCloud} a été utilisé comme outil d'analyse statique continue du
code source. Il évalue quatre dimensions : Reliability (bugs), Security
(vulnérabilités), Maintainability (code smells) et Duplication.

\begin{figure}[H]
  \centering
  \IfFileExists{figures/sonarcloud_dashboard.png}{%
    \includegraphics[width=0.88\textwidth]{figures/sonarcloud_dashboard.png}%
  }{%
    \fbox{\parbox{0.88\textwidth}{\centering\vspace{2cm}
      \textbf{[IMAGE À INSÉRER]}\\[0.4em]
      \texttt{figures/sonarcloud\_dashboard.png}\\[0.3em]
      \textit{Dashboard SonarCloud : Quality Gate Passed (vert), 0 bug, 0 vulnérabilité,
      duplication $\approx$ 2\,\%, couverture $\approx$ 83\,\%}
    \vspace{2cm}}}%
  }
  \caption{Rapport SonarCloud — Quality Gate final}
  \label{fig:sonar_dashboard}
\end{figure}

\subsection{Analyse du code}

\begin{table}[H]
  \centering
  \caption{Résultats de l'analyse SonarCloud — Quality Gate}
  \label{tab:sonar_results}
  \begin{tabularx}{\textwidth}{L{4cm}L{4cm}L{3cm}C{2cm}}
    \toprule
    \thead{Dimension} & \thead{Indicateur} & \thead{Valeur obtenue} & \thead{Statut} \\
    \midrule
    \rowcolor{lightgray}
    Reliability   & Bugs critiques        & \textbf{0}             & \cmark Passé \\
    Security      & Vulnérabilités hautes & \textbf{0}             & \cmark Passé \\
    \rowcolor{lightgray}
    Maintainability & Code smells         & \textbf{faible}        & \cmark Passé \\
    Duplication   & Taux de duplication   & \textbf{$\approx$ 2\,\%} & \cmark Passé \\
    \rowcolor{lightgray}
    Coverage      & Couverture tests      & \textbf{$\approx$ 83\,\%} & \cmark Passé \\
    \midrule
    \multicolumn{3}{l}{\textbf{Quality Gate global}} & \textbf{\cmark Passé} \\
    \bottomrule
  \end{tabularx}
\end{table}

Les corrections appliquées suite à l'analyse initiale ont porté sur : comparaisons
strictes, validation des entrées via \texttt{FormRequest}, extraction de logique
dupliquée en services partagés et suppression des logs exposant des données
sensibles.

% ============================================================
\section{DevOps}
% ============================================================

\subsection{Docker}

Le backend Laravel est conteneurisé via un \textbf{Dockerfile} dédié
(\texttt{laravel/Dockerfile}) basé sur \textbf{PHP~8.2-FPM} :

\begin{lstlisting}[style=sqlstyle,
  caption={Dockerfile Laravel (extrait)},
  label={lst:dockerfile}]
FROM php:8.2-fpm

WORKDIR /var/www

RUN apt-get update && apt-get install -y \
    git curl libpng-dev libonig-dev libxml2-dev \
    zip unzip libzip-dev mariadb-client \
  && docker-php-ext-install pdo_mysql mbstring exif pcntl bcmath gd zip

COPY --from=composer:2.7 /usr/bin/composer /usr/bin/composer

COPY --chown=www-data:www-data . /var/www

RUN composer install --no-interaction --prefer-dist \
    --optimize-autoloader --no-dev

EXPOSE 8000
CMD ["/usr/local/bin/start.sh"]
\end{lstlisting}

Le script de démarrage (\texttt{start.sh}) exécute séquentiellement
\texttt{php artisan config:cache}, \texttt{route:cache},
\texttt{migrate --force} puis \texttt{serve --host=0.0.0.0 --port=8000}.

\subsection{Docker Compose}

Pour le développement local, les trois composants de la plateforme sont orchestrés
en utilisant Docker Compose avec cinq services :

\begin{table}[H]
  \centering
  \caption{Services Docker Compose (développement)}
  \label{tab:docker_services}
  \begin{tabularx}{\textwidth}{L{3cm}L{3.5cm}X}
    \toprule
    \thead{Service} & \thead{Image} & \thead{Rôle} \\
    \midrule
    \rowcolor{lightgray}
    \texttt{laravel}  & PHP 8.2-FPM   & Backend API REST \\
    \texttt{nginx}    & Nginx 1.25    & Reverse proxy HTTP \\
    \rowcolor{lightgray}
    \texttt{mysql}    & MySQL 8       & Base de données \\
    \texttt{reverb}   & PHP 8.2-FPM   & Serveur WebSocket \\
    \rowcolor{lightgray}
    \texttt{ml}       & Python 3.11   & Microservice FastAPI \\
    \bottomrule
  \end{tabularx}
\end{table}

\subsection{Docker Hub}

Les images Docker sont publiées sur \textbf{Docker Hub} pour faciliter le déploiement
sur tout environnement de production compatible. Chaque image est versionnée par
numéro de sprint lors de la livraison.

% ============================================================
\section{CI/CD}
% ============================================================

\subsection{GitHub Actions}

Le pipeline CI/CD est configuré via \textbf{GitHub Actions} et se déclenche
automatiquement sur chaque \textit{push} et \textit{pull request} vers la branche
\texttt{main}. Il s'exécute en quatre étapes séquentielles :

\begin{table}[H]
  \centering
  \caption{Étapes du pipeline CI/CD — GitHub Actions}
  \label{tab:cicd_steps}
  \begin{tabularx}{\textwidth}{C{0.8cm}L{3.5cm}X L{2.5cm}}
    \toprule
    \thead{\#} & \thead{Étape} & \thead{Outils} & \thead{Durée moy.} \\
    \midrule
    \rowcolor{lightgray}
    1 & Lint \& style  & PHP CS Fixer, ESLint, Prettier           & $\approx$ 1\,min \\
    2 & Tests          & PestPHP (Laravel) + Vitest (React)       & $\approx$ 4\,min \\
    \rowcolor{lightgray}
    3 & Analyse qualité & SonarCloud scanner, rapport couverture  & $\approx$ 3\,min \\
    4 & Build Docker   & \texttt{docker build} + push Docker Hub  & $\approx$ 5\,min \\
    \bottomrule
  \end{tabularx}
\end{table}

\subsection{Automatisation}

Le pipeline garantit qu'aucun code ne fusionne sur \texttt{main} sans avoir
passé les trois premières étapes (lint, tests, SonarCloud). L'étape de build
Docker est déclenchée uniquement sur les pushes sur \texttt{main} (pas sur les
pull requests).

% ============================================================
\section{Déploiement}
% ============================================================

\subsection{Azure}

Le déploiement de production cible \textbf{Microsoft Azure} avec les services
suivants :

\begin{table}[H]
  \centering
  \caption{Infrastructure de déploiement — Azure}
  \label{tab:azure_infra}
  \begin{tabularx}{\textwidth}{L{4cm}L{4cm}X}
    \toprule
    \thead{Composant} & \thead{Service Azure} & \thead{Configuration} \\
    \midrule
    \rowcolor{lightgray}
    Frontend React    & Azure Static Web Apps & Build Vite, CDN mondial \\
    Backend Laravel   & Azure App Service     & Plan B2 (2 vCPU, 3,5\,Go), PHP-FPM + Nginx \\
    \rowcolor{lightgray}
    ML FastAPI        & Azure Container Apps  & 2 réplicas, Uvicorn \\
    Base de données   & Azure Database for MySQL & Flexible Server, 8\,Go RAM \\
    \rowcolor{lightgray}
    WebSocket         & Même App Service      & Laravel Reverb, port 8080 \\
    CDN / TLS         & Azure Front Door      & HTTPS, DDoS protection \\
    \bottomrule
  \end{tabularx}
\end{table}

\subsection{Architecture de production}

\begin{figure}[H]
  \centering
  \IfFileExists{figures/deploy_architecture.png}{%
    \includegraphics[width=0.90\textwidth]{figures/deploy_architecture.png}%
  }{%
    \fbox{\parbox{0.90\textwidth}{\centering\vspace{2cm}
      \textbf{[IMAGE À INSÉRER]}\\[0.4em]
      \texttt{figures/deploy\_architecture.png}\\[0.3em]
      \textit{Architecture de production Azure : utilisateur → Azure Front Door →
      Static Web Apps (React) + App Service (Laravel) + Container Apps (FastAPI)
      + Azure MySQL Flexible Server}
    \vspace{2cm}}}%
  }
  \caption{Architecture de déploiement — Production Azure}
  \label{fig:deploy_arch}
\end{figure}

% ============================================================
\section{Monitoring}
% ============================================================

\subsection{Supervision}

La supervision de production utilise les mécanismes natifs Laravel et Azure :
\begin{itemize}
  \item \textbf{Laravel Horizon} — supervision des queues (jobs de notification,
        envoi d'emails, traitement IA) ;
  \item \textbf{Azure Monitor} — métriques CPU, mémoire, latence HTTP ;
  \item \textbf{Health endpoint} — \texttt{GET /health} sur le microservice FastAPI
        retourne \texttt{\{"status": "ok"\}} pour les sondes de disponibilité.
\end{itemize}

\subsection{Journalisation}

La journalisation applicative est assurée à deux niveaux :
\begin{itemize}
  \item \textbf{Laravel Log} (Monolog) — niveaux error, warning, info ; fichiers
        rotatifs quotidiens dans \texttt{storage/logs/} ;
  \item \textbf{ActivityLog enterprise} — traçabilité complète des actions
        utilisateurs avec sévérité, IP, rôle, avant/après (décrit au Chapitre~V).
\end{itemize}

% ============================================================
\section{Résultats obtenus}
% ============================================================

\begin{table}[H]
  \centering
  \caption{Récapitulatif des métriques clés du projet}
  \label{tab:resultats}
  \begin{tabularx}{\textwidth}{L{5cm}C{3cm}C{3cm}}
    \toprule
    \thead{Indicateur} & \thead{Valeur obtenue} & \thead{Cible} \\
    \midrule
    \rowcolor{lightgray}
    User stories livrées          & \textbf{34 / 34}   & 34 \\
    Story points livrés           & \textbf{228}        & 228 \\
    \rowcolor{lightgray}
    Endpoints API REST             & \textbf{70+}        & -- \\
    Événements WebSocket           & \textbf{11}         & -- \\
    \rowcolor{lightgray}
    Tests fonctionnels             & \textbf{78 / 78}   & -- \\
    Couverture backend (SonarCloud)& \textbf{$\approx$ 83\,\%} & $\geq$ 80\,\% \\
    \rowcolor{lightgray}
    Duplication (SonarCloud)       & \textbf{$\approx$ 2\,\%}  & $\leq$ 3\,\% \\
    Bugs critiques SonarCloud      & \textbf{0}          & 0 \\
    \rowcolor{lightgray}
    Vulnérabilités hautes          & \textbf{0}          & 0 \\
    \bottomrule
  \end{tabularx}
\end{table}

% ============================================================
\section{Conclusion}
% ============================================================

\begin{chapconclusion}
Ce dernier chapitre a validé la fiabilité et la qualité de TAC-TIC IntelliPM.
Les \textbf{78 tests fonctionnels} passent tous avec succès. Le
\textbf{Quality Gate SonarCloud} est franchi au vert : 0 bug critique,
0 vulnérabilité, duplication $\approx$~2\,\%. La conteneurisation Docker,
le pipeline CI/CD GitHub Actions et l'architecture de déploiement Azure
garantissent la reproductibilité et la scalabilité de la plateforme.
\end{chapconclusion}
